\section{A generic-matrix chart in characteristic three}\label{sec:three-seed}
This section proves full monodromy in characteristic three at $d=n+2$ by making the entire Hasse--Witt matrix a set of affine coordinates on an explicit coefficient open.

Assume $n\ge3$, $d=n+2$ and $p=3$. The positive indices are $a_i=\mathbf1+e_i$, $0\le i\le n$, and $h=n+1$. The convention of \eqref{eq:pn-1-4} gives
\begin{equation}\label{eq:three-matrix}
B_3(f)_{i,j}=[x^{M_{i,j}}]f^2,
\qquad M_{i,j}=2\mathbf1+3e_j-e_i.
\end{equation}
For $i=j$, the target exponent has entry $4$ at $j$ and entry $2$ elsewhere. For $i\ne j$, it has entry $5$ at $j$, entry $1$ at $i$, and entry $2$ elsewhere. In particular it has at most one entry at least four.

Call a degree-$d$ monomial \emph{low} if every exponent is at most three, and \emph{high} otherwise. Write $f=g+z$ accordingly. Two high monomials have either an exponent sum at least eight in one coordinate, or exponents at least four in two coordinates. Their product can never have exponent $M_{i,j}$. Consequently
\begin{equation}\label{eq:three-linear}
B_3(g+z)=B_3(g)+2L_g(z),
\qquad L_g(z)_{i,j}=[x^{M_{i,j}}]gz.
\end{equation}
A high monomial with exponent at least four in position $j$ can contribute only to column $j$. A monomial with two such positions contributes to no column.

\subsection{Explicit invertible coefficient minors}
Fix $j$ and choose three distinct other indices $a,b,c$. Let $R$ be the remaining set, of size $n-3$, and put $P_R=\prod_{r\in R}x_r$. Use the low polynomial
\[
g_j=(x_ax_b^2x_c^2+x_jx_a^2x_c^2+x_jx_a^2x_b^2)P_R.
\]
It has degree $n+2$ and all its exponents are at most two. Select the following $h$ high monomials:
\begin{equation}\label{eq:three-monomials}
\begin{aligned}
Z_0&=x_j^5P_R,&Z_a&=x_j^4x_aP_R,\\
Z_b&=x_j^4x_bP_R,&Z_c&=x_j^4x_cP_R,\\
Z_r&=x_j^5x_a\prod_{s\in R\setminus\{r\}}x_s
 &&(r\in R).
\end{aligned}
\end{equation}
They are distinct degree-$n+2$ monomials, each with precisely one exponent at least four, in position $j$. Ordering rows as $j,a,b,c,R$ and columns as $0,a,b,c,R$, the selected matrix of $L_{g_j}$ in column $j$ is
\begin{equation}\label{eq:three-minor}
\begin{pmatrix}
0&1&0&0\\1&0&0&0\\0&0&1&0\\0&0&0&1
\end{pmatrix}\oplus I_{n-3}.
\end{equation}
Write $\nu(Z)$ for the exponent vector of a monomial $Z$. Here is the complete coefficient check. For the first four rows and columns, each variable in $R$ has target exponent two and selected exponent one, so the remaining exponent is one. The calculation reduces to the four variables $j,a,b,c$. Its four nonzero differences are
\[
\begin{aligned}
M_{j,j}-\nu Z_a&=\nu(x_ax_b^2x_c^2P_R),&
M_{a,j}-\nu Z_0&=\nu(x_ax_b^2x_c^2P_R),\\
M_{b,j}-\nu Z_b&=\nu(x_jx_a^2x_c^2P_R),&
M_{c,j}-\nu Z_c&=\nu(x_jx_a^2x_b^2P_R),
\end{aligned}
\]
Every other difference has a negative exponent or is absent from $g_j$. For row $r\in R$ and column $Z_r$, the difference is the first monomial of $g_j$. For a different row and column $Z_r$, coordinate $r$ would require exponent two in $g_j$, whereas its exponent there is one. Conversely row $r$ and a first-four column would require exponent zero there. This proves \eqref{eq:three-minor}; its determinant is $-1$.

For the general low polynomial, let $D_j$ be this minor determinant. It is a nonzero polynomial. The product $\prod_jD_j$ is nonzero in the integral low-coefficient polynomial ring, so there is a nonempty open $B_{\mathrm{low}}$ on which every column minor is invertible. The special polynomials $g_j$ used to verify the separate minors need not be the same polynomial.

\subsection{The coefficient chart and its connected Lang cover}
The chosen monomials are distinct between columns, because their unique exponent at least four occurs at different positions. Use their $h^2$ coefficients as distinguished high coefficients and put the remaining high coefficients in an affine space $B_{\mathrm{high}}$. Set $B=B_{\mathrm{low}}\times B_{\mathrm{high}}$.

Equation~\eqref{eq:three-linear} gives an isomorphism
\begin{equation}\label{eq:three-chart}
B\times\mathbf A^{h^2}\xrightarrow{\ \sim\ }B\times\operatorname{Mat}_h,
\qquad f\longmapsto (\text{remaining coefficients},B_3(f)).
\end{equation}
Its linear part in the distinguished coefficients is block diagonal by columns, with blocks twice the invertible minors $D_j$. The inverse solves those linear systems, using only the determinants already inverted. Other high coefficients only change the known affine constant term. Thus \eqref{eq:three-chart} is a regular coordinate isomorphism, not just a dominant map. It describes an open of the full affine coefficient space; no coefficient has been fixed to zero.

The ordinary part is $B\times\operatorname{GL}_h$. Its Lang basis cover, by the elementary fixed-vector calculation of Section~\ref{sec:pn-2}, has total space $B\times\operatorname{GL}_h$, with map
\[
(b,T)\longmapsto\bigl(b,T(T^{[3]})^{-1}\bigr).
\]
It is a finite \'etale $\operatorname{GL}_h(\mathbf F_3)$-torsor and its total space is geometrically integral. Hence its geometric image is the full group.

The ordinary coefficient space is integral and normal. Its smooth hypersurface open is nonempty by Section~\ref{sec:pn-3}, and it meets the nonempty coefficient chart above. Restricting the connected finite \'etale basis cover to this nonempty open preserves connectedness. Passing from nonzero affine equations to projective equations places this full affine image inside the projective image. The latter is therefore also $\operatorname{GL}_h(\mathbf F_3)$. The arithmetic image contains it and is bounded by the same group. This proves the degree-$d=n+2$ characteristic-three case of Theorem~\ref{thm:higher-images}.
